\documentclass[../main.tex]{subfiles}
\begin{document}

\chapter{ストレージ}
\lessoninfo{
  video={Lesson 16，動画 1--14},
  textbook={H\&P \cite{hennessy2017} §2.2，§D.1--D.2},
  keywords={磁気ディスク，プラッタ，シリンダ，セクタ，シーク時間，回転待ち時間，光ディスク，磁気テープ，順次アクセス，フラッシュメモリ，SSD，ハイブリッドディスク，I/O バス，メザニンバス},
}

第15章では，主記憶が DRAM で作られており，電源を切れば直ちに内容を失う
ことを見た．残しておかなければならないプログラムとデータは，ストレージ
装置に保存しなければならない．ストレージ装置は，仮想記憶（第13章）が物理メモリ
に収めきれないページも保持する．この章では，主要なストレージ技術---磁気
ディスク，光ディスク，磁気テープ，および DRAM やフラッシュメモリによる
半導体ストレージ---を，そのアクセス時間を何が決めるか，そしてそれがどの
ように進歩するかに重点を置いて述べ，ストレージとその他の I/O デバイスが
プロセッサにどのように接続されるかを示す．

\section{ストレージ}
\label{sec:l16-storage}

ストレージはデータを\emph{永続的に}保持する．主記憶とは異なり，ストレージ
装置は計算機の電源を切っても内容を保つ．\source{Lesson 16，動画 1--2；
H\&P §D.1．}その用途は2つある．
\begin{description}
  \item[ファイル] プログラム，それが扱うデータ，システムとアプリケーション
    の設定はファイルに保存され，これらはあるセッションから次のセッションへ
    残らなければならない．
  \item[仮想記憶] 物理メモリに収まらないページはストレージに保持され，
    ページフォールトはそのページをメモリへ読み戻す（第13章）．
\end{description}

ストレージ装置についてアーキテクトが注目する性質は2つある．その\emph{性能}は
第2章の2つの指標で記述される．すなわち，毎秒処理されるデータ量あるいは
要求数であるスループットと，1回のアクセスに要する時間であるレイテンシ
である．ストレージのスループットは時代とともに向上するが，プロセッサの
速度ほど速くは向上しない．ストレージのレイテンシの改善はさらに遅く，DRAM
のレイテンシの改善よりもなお遅い．そのため，プロセッサとストレージの隔たり
は第1章のメモリウォールよりもさらに大きい．\margin{読み出しレイテンシ：
DRAM は 60--\SI{90}{\nano\second}，フラッシュ SSD はおよそ
\SI{100}{\micro\second}，ディスクはおよそ \SI{10}{\milli\second} である
（7-cpu.com，ベンダのデータシート）．}その\emph{信頼性}も少なくとも
同じくらい重要である．ストレージはデータが安全に保たれると期待される場所で
あり，ファイルを失うことは，プログラムが落ちることよりも深刻な場合が多い．
信頼性は第17章の主題である．

\section{磁気ディスク}
\label{sec:l16-disk}

\term[じきでぃすく]{磁気ディスク}[magnetic disk]%
\index{はーどでぃすく@ハードディスク|see{磁気ディスク}}（ハードディスク）
は，磁性体の被膜の微小な領域の磁化の向きとしてビットを記録する．
\source{Lesson 16，動画 3；H\&P §D.1．}この被膜で覆われた回転する円板と，
その上を動くヘッドとから作られる（\cref{fig:l16-disk}）．

\begin{figure}[htbp]
  \centering
  % Structure of a magnetic disk.  (a) Edge-on view: platters on a spindle,
% one head per surface on arms that form the head assembly; the dashed
% lines mark one cylinder.  (b) Top view of one surface: concentric tracks
% divided into sectors; one track and one sector are highlighted.
\begin{tikzpicture}[x=1cm,y=1cm, font=\footnotesize\sffamily, >={Stealth[length=1.6mm]},
  lab/.style={font=\scriptsize\sffamily, text=ln-muted},
  pin/.style={draw=ln-muted, thin}]
  % ================================================= (a) edge-on view
  \node[anchor=west] at (-0.2,3.35) {\iflangja{(a) 側面}{(a) side view}};
  % spindle
  \draw[fill=ln-muted!35] (1.45,-0.15) rectangle (1.75,2.45);
  % platters
  \foreach \y in {0.4,1.2,2.0} {
    \draw[fill=ln-accent!30] (0.1,\y-0.05) rectangle (3.1,\y+0.05);
  }
  % one cylinder: all tracks at the same radius
  \draw[ln-accent, dashed, thick] (2.35,0.15) -- (2.35,2.25);
  \draw[ln-accent, dashed, thick] (0.85,0.15) -- (0.85,2.25);
  % head assembly: actuator and arms, one head per surface
  \draw[fill=ln-box] (4.3,0.12) rectangle (4.55,2.28);
  \foreach \y in {0.4,1.2,2.0} {
    \foreach \s in {1,-1} {
      \draw[very thick, ln-muted] (4.3,\y+\s*0.15) -- (2.5,\y+\s*0.15);
      \fill[ln-accent] (2.22,\y+\s*0.07) rectangle (2.5,\y+\s*0.19);
    }
  }
  % rotation
  \draw[->, thick] (1.05,2.62) arc[start angle=190, end angle=350, x radius=0.55, y radius=0.14];
  % labels
  \node[lab, anchor=south] at (1.6,2.75) {\iflangja{スピンドル}{spindle}};
  \draw[pin] (0.35,0.35) -- (0.35,-0.25);
  \node[lab, anchor=north] at (0.35,-0.22) {\iflangja{プラッタ}{platter}};
  \draw[pin] (2.35,0.15) -- (2.35,-0.25);
  \node[lab, anchor=north] at (2.35,-0.22) {\iflangja{シリンダ}{cylinder}};
  \draw[pin] (2.42,2.2) -- (3.2,2.74);
  \node[lab, anchor=south] at (3.3,2.72) {\iflangja{ヘッド}{head}};
  \node[lab, anchor=north] at (4.43,0.05) {\iflangja{ヘッドアセンブリ}{head assembly}};
  % ================================================= (b) top view
  \begin{scope}[xshift=6.3cm]
    \node[anchor=west] at (-0.2,3.35) {\iflangja{(b) 1 面の上面}{(b) top view of one surface}};
    \coordinate (c) at (2.0,1.2);
    \draw[fill=ln-box] (c) circle[radius=1.5];
    \foreach \r in {0.5,0.65,...,1.4} {
      \draw[ln-rule] (c) circle[radius=\r];
    }
    \foreach \a in {0,30,...,330} {
      \draw[ln-rule] ($(c)+(\a:0.5)$) -- ($(c)+(\a:1.4)$);
    }
    % one track and one sector
    \draw[ln-accent, thick] (c) circle[radius=0.95];
    \fill[ln-accent] ($(c)+(30:1.02)$) arc[start angle=30, end angle=60, radius=1.02]
      -- ($(c)+(60:0.88)$) arc[start angle=60, end angle=30, radius=0.88] -- cycle;
    % spindle
    \draw[fill=ln-muted!35] (c) circle[radius=0.3];
    % head on its arm
    \draw[very thick, ln-muted] (3.75,-0.05) -- ($(c)+(-12:0.95)$);
    \fill[ln-muted] (3.75,-0.05) circle[radius=0.08];
    \fill[ln-accent] ($(c)+(-12:0.95)$) circle[radius=0.07];
    % rotation
    \draw[->, thick] ($(c)+(105:1.65)$) arc[start angle=105, end angle=165, radius=1.65];
    % labels
    \draw[pin] ($(c)+(45:0.95)$) -- (3.45,2.7);
    \node[lab, anchor=south] at (3.45,2.67) {\iflangja{セクタ}{sector}};
    \draw[pin] ($(c)+(215:0.95)$) -- (0.1,-0.2);
    \node[lab, anchor=north] at (0.1,-0.17) {\iflangja{トラック}{track}};
    \draw[pin] ($(c)+(-12:0.95)$) -- (3.55,1.45);
    \node[lab, anchor=west] at (3.52,1.5) {\iflangja{ヘッド}{head}};
  \end{scope}
\end{tikzpicture}

  \caption{磁気ディスクの構造．(a)~スピンドル上のプラッタと，1面につき
    1つのヘッド．すべてのヘッドのアームはヘッドアセンブリをなし，破線は
    1つのシリンダを示す．(b)~1つの面．同心円状のトラックに分けられ，各
    トラックはセクタに分けられる．1つのトラックとそのセクタの1つを強調して
    ある．}
  \label{fig:l16-disk}
\end{figure}

被膜は，剛体の円板である\term[ぷらった]{プラッタ}[platter]の両面を覆う．
複数のプラッタが共通の\term[すぴんどる]{スピンドル}[spindle]に積み重ねられ，
スピンドルはモータによって一定の速度で回されるので，すべてのプラッタは
一体となって回転する（\cref{fig:l16-disk}a）．

各面は，それぞれ専用の\term[じきへっど]{磁気ヘッド}[disk head]によって読み
書きされる．ヘッドはアームの先端に取り付けられ，面に触れることなくごく
わずかな距離だけ浮いている．回転している面にヘッドが触れれば，面を傷つけて
しまう．すべてのアームは1つのアクチュエータに取り付けられ，\emph{ヘッド
アセンブリ}として一体で動いて，プラッタの外周からスピンドルまでの任意の
位置にヘッドを置くことができる．

ヘッドアセンブリを1つの位置に留めておくと，プラッタが回転するにつれて各
ヘッドは自身の面上の1つの円周をたどる．この円周が
\term[とらっく]{トラック}[track]であり，1つの面は多数の同心円状のトラックを
持つ（\cref{fig:l16-disk}b）．ヘッドは一体で動くので，各面に1つずつ，常に
スピンドルから等しい距離にあるトラックの上に位置する．これらのトラックは
まとめて\term[しりんだ]{シリンダ}[cylinder]をなす．すなわち，ヘッド
アセンブリを動かすことがシリンダの選択であり，ヘッドの1つを選ぶことが面の
選択，したがってそのシリンダ内の1本のトラックの選択である．

トラックに沿って，データは固定長のブロックに分割される．このブロックを
\term[せくた]{セクタ}[sector]といい，ディスクが読み書きする最小の単位で
ある．セクタは短いプリアンブルで始まり，ディスクはこれによってセクタの
先頭を認識し，そのビットに同期する．続いてデータが並び，セクタの末尾には
誤り訂正符号が置かれる．ディスクはこれを用いて，読み出したデータを検査し，
可能であれば訂正する（第17章）．

ディスクの容量は，この構成から直ちに求まる．
\begin{equation}
  \text{容量} \;=\; \text{プラッタ数} \times 2
    \times \frac{\text{トラック数}}{\text{面}}
    \times \frac{\text{セクタ数}}{\text{トラック}}
    \times \frac{\text{バイト数}}{\text{セクタ}} ,
  \label{eq:l16-capacity}
\end{equation}
ここで係数 2 はプラッタの両面を数えたものであり，1面当たりのトラック数は
シリンダ数に等しい．\margin{実際のディスクは，長い外周のトラックにより
多くのセクタを置く．その場合，1トラック当たりのセクタ数は平均値である．}

\begin{example}
  \label{ex:l16-capacity}
  あるディスクが2枚のプラッタを持ち，1面当たり \num{100000} トラック，
  1トラック当たり平均 \num{2000} セクタ，1セクタ当たり \num{4096} バイト
  であるとする．\cref{eq:l16-capacity} により $2\times 2\times
  \num{100000}\times \num{2000} = \num{8e8}$ セクタ，すなわち
  $\num{8e8}\times \num{4096} \approx \num{3.3e12}$ バイト，およそ
  3.3~TB を保持する．\caution{ディスクの容量は10進で
  ある．1~TB は $2^{40}$ ではなく $10^{12}$ バイトなので，4~TB のドライブが
  保持するのは 3.64~TiB にすぎない．}
\end{example}

\section{ディスクのアクセス時間}
\label{sec:l16-access}

セクタを読み書きするには，ディスクは正しいヘッドを回転する面の正しい位置に
持っていかなければならず，その後データが主記憶まで運ばれなければならない．
\source{Lesson 16，動画 4--5；H\&P §D.1．}プラッタがすでに回転している
ディスクでは，1回のアクセスの時間は5つの部分からなる．
\begin{description}
  \item[シーク時間] \term[しーくじかん]{シーク時間}[seek time]とは，
    ヘッドアセンブリが，目的のセクタを含むシリンダまでヘッドを動かすのに
    要する時間である．ヘッドが移動する距離とともに増える．
  \item[回転待ち時間] \term[かいてんまちじかん]{回転待ち時間}%
    [rotational latency]とは，ヘッドがトラックに到達してから，セクタの
    先頭がヘッドの下に回ってくるまでの時間である．
  \item[データ読み出し時間] セクタは回転の速度でヘッドの下を通過し，その
    間に読み出される（あるいは書き込まれる）．連続する複数のセクタへの
    アクセスは，それに応じて長くかかる．
  \item[コントローラ時間] ディスクコントローラが要求を処理する時間で
    ある．たとえば読み出した各セクタの誤り訂正符号の検査がこれに当たる．
  \item[I/O バス時間] データはディスクと主記憶の間を I/O バス
    （\cref{sec:l16-io}）を通じて転送される．
\end{description}
\begin{equation}
  t_{\text{access}} \;=\; t_{\text{seek}} + t_{\text{rot}} + t_{\text{read}}
    + t_{\text{controller}} + t_{\text{bus}} .
  \label{eq:l16-access}
\end{equation}
ディスクが先行する要求でまだふさがっていれば，新しい要求は自身のアクセスが
始まる前に待ち行列で待つことにもなる．この待ち行列遅延---先行するアクセスが
終わり，ヘッドが再び動けるようになるまでの時間---は無視できないことが多い．

回転に依存する2つの項は回転速度から求まる．回転速度は通常，毎分回転数
（RPM）で与えられる．1回転は $60/\text{RPM}$ 秒かかる．\tip{ミリ秒では
1回転は $\num{60000}/\text{RPM}$ である．\num{7200}~RPM なら
\SI{8.33}{\milli\second} であり，平均の回転待ち時間はおよそ
\SI{4.2}{\milli\second} となる．}ヘッドが到達したとき，セクタの先頭は
トラック上のどこにあるか分からないので，回転待ち時間は 0 から1回転の間にあり，
平均では半回転である．連続する $n$ セクタの読み出しには，1回転の
$n/(\text{トラック当たりのセクタ数})$ の割合だけかかる．
\begin{equation}
  \overline{t_{\text{rot}}} \;=\; \frac{1}{2}\cdot\frac{60}{\text{RPM}} ,
  \qquad
  t_{\text{read}} \;=\; \frac{n}{\text{トラック当たりのセクタ数}}\cdot
    \frac{60}{\text{RPM}} .
  \label{eq:l16-rot}
\end{equation}

\begin{example}
  \label{ex:l16-access}
  あるディスクが \num{10000}~RPM で回転し，平均シーク時間が
  \SI{4}{\milli\second} であるとする．ある要求が，500 セクタを持つ
  トラックの連続する 100 セクタを読む．コントローラは
  \SI{0.3}{\milli\second}，I/O バスでの転送は \SI{0.5}{\milli\second} を
  要する．1回転は $\num{60000}/\num{10000} = \SI{6}{\milli\second}$
  かかるので，\cref{eq:l16-rot} により平均の回転待ち時間は
  \SI{3}{\milli\second}，データ読み出し時間は $100/500 \times
  \SI{6}{\milli\second} = \SI{1.2}{\milli\second}$ である．
  \cref{eq:l16-access} により，このアクセスには平均で $4 + 3 + 1.2 + 0.3
  + 0.5 = \SI{9}{\milli\second}$ かかる（\cref{fig:l16-access-time}）．
  もし 100 セクタがディスク上に散在していれば，そのそれぞれに固有のシークと
  回転待ち時間が必要となり，この要求には \SI{700}{\milli\second} を超える
  時間がかかったはずである．\margin{ドライブは一度に1つの
  アクセスしか処理しない．\SI{9}{\milli\second} は毎秒およそ 110 要求で
  ある．ベンダはこれを IOPS と呼ぶ．フラッシュ SSD は \num{100000} 以上に
  達する．}
\end{example}

\begin{figure}[htbp]
  \centering
  % Components of one disk access, drawn to scale for the numbers of
% Example ex:l16-access (seek 4 ms, rotational latency 3 ms, data read
% 1.2 ms, controller 0.3 ms, I/O bus 0.5 ms).  1 ms = 1.2 cm.
\begin{tikzpicture}[x=1.2cm,y=1cm, font=\footnotesize\sffamily,
  lab/.style={font=\scriptsize\sffamily, text=ln-muted},
  pin/.style={draw=ln-muted, thin},
  seg/.style={draw, minimum height=0.55cm}]
  % segments (x in ms)
  \draw[fill=ln-accent!40] (0,0) rectangle (4,0.55);
  \draw[fill=ln-accent!20] (4,0) rectangle (7,0.55);
  \draw[fill=ln-accent!10] (7,0) rectangle (8.2,0.55);
  \draw[fill=ln-box] (8.2,0) rectangle (8.5,0.55);
  \draw[fill=white] (8.5,0) rectangle (9,0.55);
  \node at (2,0.275) {\iflangja{シーク時間}{seek time}};
  \node at (5.5,0.275) {\iflangja{回転待ち時間}{rotational latency}};
  \node[lab, anchor=south] at (7.6,0.6) {\iflangja{データ読み出し}{data read}};
  \draw[pin] (8.35,0.55) -- (8.35,1.0);
  \node[lab, anchor=south] at (8.35,0.97) {\iflangja{コントローラ}{controller}};
  \draw[pin] (8.75,0.55) -- (8.95,0.72);
  \node[lab, anchor=south west] at (8.85,0.62) {\iflangja{I/O バス}{I/O bus}};
  % time axis
  \draw (0,-0.1) -- (9,-0.1);
  \foreach \t in {0,...,9} {
    \draw (\t,-0.1) -- (\t,-0.18);
    \node[lab, anchor=north] at (\t,-0.16) {\t};
  }
  \node[lab, anchor=north west] at (9.1,-0.16) {ms};
  % mechanical vs electronic parts
  \draw[ln-muted] (0,-0.6) -- (0,-0.68) -- (8.15,-0.68) -- (8.15,-0.6);
  \node[lab, anchor=north] at (4.1,-0.7) {\iflangja{機械的な動作（ヘッドとプラッタ）}{mechanical (head and platter movement)}};
  \draw[ln-muted] (8.25,-0.6) -- (8.25,-0.68) -- (9,-0.68) -- (9,-0.6);
  \node[lab, anchor=north] at (8.61,-0.7) {\iflangja{電子的}{electronic}};
\end{tikzpicture}

  \caption{\cref{ex:l16-access} のアクセスを実寸で描いたもの．
    \SI{9}{\milli\second} のうち \SI{8.2}{\milli\second} を機械的な動作が
    占める．}
  \label{fig:l16-access-time}
\end{figure}

\tip{\SI{1}{\milli\second} は，\SI{1}{\giga\hertz} 当たり100万サイクルと
読めばよい．上の \SI{9}{\milli\second} は \SI{3}{\giga\hertz} の
プロセッサで2700万サイクルである．}

\begin{keypoint}
  ディスクのアクセスは機械的な動作---ヘッドアセンブリの移動と，プラッタが
  回ってくるのを待つこと---に支配される．それはミリ秒，すなわちプロセッサ
  のクロックで数百万サイクルを要する．ページフォールト（第13章）がこれほど
  高価なのはこのためである．
\end{keypoint}

\section{磁気ディスクの動向}
\label{sec:l16-trends}

磁気ディスクの構成要素は，きわめて異なる速度で改善する．\source{Lesson
16，動画 6；H\&P §D.2．}\sidenote{磁気ディスクは，今日
なお使われている計算機の構成要素のうち最も古い．IBM の 350 ディスク記憶装置
（1956 年）は，24 インチのプラッタ 50 枚に約 5~MB を蓄え，\num{1200}~RPM で
回した．ヘッドは1つだけで，トラックに到達するのにおよそ \SI{0.6}{\second} を
要した．容量はその後6桁伸びたが，アクセス時間は2桁も縮まっていない．}
\begin{description}
  \item[容量] は，面上にどれだけ密にビットを記録するかで決まる．長年に
    わたり1〜2年でおよそ倍になってきたが，2010 年頃からこの成長はかなり
    鈍っている．
  \item[シーク時間] は通常 \SIrange[range-phrase={〜}]{5}{10}{\milli\second} であり，ごく
    ゆっくりとしか改善しない．\caution{公称のシーク時間はランダムなシークの
    平均である．隣接するトラックへのシークは \SI{1}{\milli\second} 未満，
    全ストロークのシークは平均のおよそ2倍かかる．}その改善分は主に，
    プラッタの直径を小さくして
    ヘッドの移動距離を短くすることから来る．
  \item[回転速度] はゆっくりと上がり，およそ \num{5000}~RPM から，最も
    速いドライブでは \num{10000}--\num{15000}~RPM になった．大半の
    ドライブは \num{5400}--\num{7200}~RPM で回る．回転を速くするには
    より強い材料が必要であり，騒音と発熱と消費電力も増える．
  \item[コントローラ時間とバス時間] は電子的であり，順調に改善するが，
    アクセス時間のうちのわずかな部分でしかない
    （\cref{fig:l16-access-time}）．
\end{description}

したがってディスクのアクセスのレイテンシは，Moore の法則（第1章）が述べる
トランジスタ技術の進歩ではなく，機械と材料によって決まる．容量が長年に
わたり急速に伸びたにもかかわらず，レイテンシはプロセッサの速度よりはるかに
ゆっくりとしか改善しない．
記録密度を高めれば，1回転でヘッドの下を通る
バイト数が増えるので，一定量のデータの読み出し時間は確かに短くなる．
しかし，シーク時間にも回転待ち時間にも影響しない．

\section{光ディスク}
\label{sec:l16-optical}

\term[ひかりでぃすく]{光ディスク}[optical disk]は，磁気ディスクのプラッタ
1枚のように回転する1枚の円板であり，データはトラックに沿って並ぶ．ただし
ビットは，面が光を反射する仕方を変えるマークとして記録され，磁気ヘッドでは
なくレーザで読み出される．\source{Lesson 16，動画 7．}レーザはディスクから
離れた位置で働き，透明な保護層を通して焦点を結ぶので，面に付いた埃や指紋の
害は，ドライブ内に密閉されていなければならない磁気ディスクのプラッタの場合
よりはるかに小さい．そのため光ディスクはドライブから取り出し，別のドライブ
へ持ち運べる．すなわち\emph{可搬}である．\margin{CD，DVD，Blu-ray ディスク
では，データは1本のらせん状のトラックに沿って並ぶ．}

可搬性には代償がある．あるドライブで書かれたディスクは他のすべてのドライブ
で読めなければならないので，記録形式は業界標準---CD，次いで DVD，次いで
Blu-ray---によって固定されていなければならない．\margin{CD（1982 年）は
700~MB，DVD（1996 年）は 4.7~GB，Blu-ray（2006 年）は1層あたり 25~GB で
ある．標準どうしの間隔はおよそ10年，容量の差はおよそ6倍である．}
より高密度な記録技術が役に
立つのは，新しい標準が合意され，それに対応するドライブと媒体とコンテンツが
普及した後に限られ，それには何年もかかる．したがって光ディスクの進歩は，
長い間隔をおいた大きな飛躍という形をとる．磁気ディスクのプラッタはドライブ
から出ることがない．製造者は，ドライブが同じ標準インタフェースで接続される
限り，次の製品でより優れた記録技術を採用できる．

\section{磁気テープ}
\label{sec:l16-tape}

\term[じきてーぷ]{磁気テープ}[magnetic tape]は，磁性体を塗布した細長い帯に
データを記録する．帯はカートリッジに巻かれて収められ，ドライブの読み書き
ヘッドの前を引き出される．\source{Lesson 16，動画 8--9．}テープは容量が
大きく，カートリッジは交換できる．1台のドライブで，安価なカートリッジを
何本でも読み書きできる．その欠点はデータへの到達の仕方にある．

\begin{definition}[順次アクセス]
  データに順番にしか到達できないとき，その媒体は
  \term[じゅんじあくせす]{順次アクセス}[sequential access]を提供すると
  いう．ある位置にアクセスするには，現在の位置と要求された位置の間にある
  データすべてを通り過ぎるまで媒体を動かさなければならず，アクセス時間は
  その距離とともに増える．
\end{definition}

対照的にディスクは，間にあるデータを通り過ぎることなく，任意のセクタに直接
到達する．1回のシークと高々1回転は，セクタがどこにあっても高々数十ミリ秒で
ある（\cref{ex:l16-tape}）．それでもアクセス時間はシーク距離に依存するので，
ディスクがランダムアクセス（第15章）を提供するのは近似的にすぎないが，
テープよりははるかにそれに近い．順次アクセスが特によく合う用途が1つある．
\emph{バックアップ}である．バックアップでは大量のデータが一度に通して
書かれ，読み戻されることは稀で，読み戻すときもたいていは一度に通して読む．
\margin{LTO-9 のカートリッジ（2021 年）は 18~TB を保持し，400~MB/s で流す
ので，端から端まで
読むのにおよそ12時間かかる（LTO コンソーシアム）．}
順次アクセスは仮想記憶には適さない．仮想記憶のページフォールトは任意の位置
からページを読むからである．したがってテープは，プログラムやファイルを使う
ためのストレージではなく，バックアップ用のストレージとして使われる．

\begin{example}
  \label{ex:l16-tape}
  \SI{1000}{\metre} のテープが
  \SI[per-mode=symbol]{5}{\metre\per\second} で巻き取られるとする．
  テープが先頭にあるとき，中央のデータには \SI{100}{\second} 後，末尾の
  データには \SI{200}{\second} 後に到達する．最長のシークが
  \SI{15}{\milli\second}，1回転が \SI{6}{\milli\second} のディスクは，
  どのセクタにもおよそ \SI{21}{\milli\second} 以内で到達する．これは
  テープが末尾に到達するよりおよそ \num{9500} 倍速い．バックアップ全体を
  先頭から末尾まで読む場合には，巻き取りの時間は問題にならない．
\end{example}

テープは次第に置き換えられつつある．生産量が少ないため，ギガバイト当たりの
コストは，膨大な数が作られる磁気ディスクほど速くは下がらない．ディスク---
たとえば USB で接続する外付けディスク---は，バックアップの役割の多くを
引き受けられるほど安価になった．

\section{SSD}
\label{sec:l16-ssd}

半導体メモリには可動部分がなく，そのアクセスには，ディスクのミリ秒ではなく，
ナノ秒からミリ秒に満たない程度の時間しかかからない．\source{Lesson 16，
動画 10；H\&P §2.2，§D.1．}しかしギガバイト当たりのコストははるかに高い．
DRAM のギガバイト当たりの価格は磁気ディスクの百倍程度であり，一方で
レイテンシはおよそ \num{100000} 分の1である（数十ナノ秒に対してミリ秒）．
半導体メモリで作られたストレージは，速度がこのコストに見合う場所で使われる．

\begin{definition}[SSD]
  \term{SSD}[solid-state disk]とは，可動部分を持たず半導体メモリで作られた
  ストレージ装置であって，ディスクのように使われるものである．データを
  ブロック単位で保存し，SATA のようなディスクインタフェース，あるいは
  PCI Express を通じて直接接続される．
\end{definition}

SSD は2種類のメモリのいずれでも作れる．
\begin{description}
  \item[電池付きの DRAM] DRAM は揮発性なので，計算機の電源が切れている間は
    電池が電力を供給し続ける．この種の SSD はきわめて高速で，ギガバイト
    当たりの価格もきわめて高い．データが安全なのは電池が保つ間だけであり，
    長期間データを保つ用途には向かない．
  \item[フラッシュメモリ] 大半の SSD は
    \term[ふらっしゅめもり]{フラッシュメモリ}[flash memory]を使う．これは
    電源がなくても内容を保つ半導体メモリであり，ビットはトランジスタに
    閉じ込められた電荷として蓄えられ，そこに何年も留まる．
\end{description}

\caution{フラッシュメモリの「フラッシュ」は，パイプラインのフラッシュ
（第3章）とは何の関係もない．}

フラッシュの SSD は DRAM より遅いが，それでも磁気ディスクよりはるかに速く，
消費電力が小さく，電池なしでデータを保つ．\margin{フラッシュは上書きが
できない．数 MB のブロック全体をまず消去しなければならず，1つのブロックが
耐えられる消去は数千回程度である．}ギガバイト当たりのコストは DRAM
とディスクの中間にあるので，同じ価格ではディスクよりずっと小さい容量しか
得られない．長い間，典型的なフラッシュ SSD の容量はギガバイト単位であり，
ディスクはテラバイト単位であった．

\section{磁気ディスクとフラッシュのハイブリッド}
\label{sec:l16-hybrid}

磁気ディスクとフラッシュメモリの長所は互いに補い合う．\source{Lesson 16，
動画 11--12．}磁気ディスクは，ギガバイト当たりの低いコストで巨大な容量を
提供するが，遅く，プラッタが回っている間はかなりの電力を消費し，動作中の
衝撃に弱い．フラッシュメモリは速く，消費電力が小さく，可動部分を持たないが，
ギガバイト当たりのコストは高い．

\begin{definition}[ハイブリッドディスク]
  \term[はいぶりっどでぃすく]{ハイブリッドディスク}[hybrid disk]とは，
  磁気ディスクと，それより少量のフラッシュメモリとを1つの装置にまとめた
  ものである．磁気ディスクが全データを保持し，フラッシュメモリは最も頻繁に
  使われるデータのキャッシュとして働く（\cref{fig:l16-hybrid}）．
\end{definition}

時間的局所性と空間的局所性（第12章）は，メモリアクセスと同様にファイルにも
当てはまる．オペレーティングシステム，日常的に使うアプリケーション，作業中
のファイルが，アクセスの大半を占める．これらをフラッシュに置いておけば，
大半の要求はフラッシュの速度で，機械的な動作なしに処理される一方，容量と
ギガバイト当たりのコストはディスクのそれに近いままである．要求がフラッシュ
から処理されている間は，ディスクの回転を止めることさえでき，電力を節約し，
装置は衝撃に強くなる．そのときディスクを必要とする要求が来れば，プラッタが
再び回り出すまで，しばしば数秒待つことになる．

\begin{marginfigure}
  % Hybrid disk: a controller that serves most requests from a flash cache
% holding the frequently used data; the magnetic disk holds all data.
\resizebox{\linewidth}{!}{%
\begin{tikzpicture}[x=1cm,y=1cm, font=\footnotesize\sffamily, >={Stealth[length=1.6mm]},
  box/.style={draw, fill=ln-box, align=center, inner sep=3pt, minimum height=0.9cm},
  lab/.style={font=\scriptsize\sffamily, text=ln-muted, align=center}]
  \node[box, minimum width=4.2cm] (ctl) at (0,0) {\iflangja{コントローラ}{controller}};
  \node[box, fill=ln-accent!20, minimum width=1.9cm] (flash) at (-1.15,-1.9)
    {\iflangja{フラッシュ}{flash}\\[-1pt]{\scriptsize\iflangja{よく使うデータ}{frequently used}}};
  \node[box, minimum width=1.9cm] (disk) at (1.15,-1.9)
    {\iflangja{磁気ディスク}{magnetic disk}\\[-1pt]{\scriptsize\iflangja{全データ}{all data}}};
  \draw[<->, very thick, ln-accent] (-1.15,-0.45) -- (flash.north)
    node[lab, midway, left] {\iflangja{大部分}{most}};
  \draw[<->] (1.15,-0.45) -- (disk.north)
    node[lab, midway, right] {\iflangja{残り}{rest}};
  \draw[<->, thick] (0,1.2) -- (ctl.north) node[lab, pos=0.25, right] {\iflangja{I/O バス}{I/O bus}};
  \node[draw=ln-muted, dashed, rounded corners=2pt, fit=(ctl)(flash)(disk), inner sep=5pt] (hyb) {};
  \node[lab, anchor=north] at (hyb.south) {\iflangja{ハイブリッドディスク}{hybrid disk}};
\end{tikzpicture}}

  \caption{ハイブリッドディスク．大半の要求はフラッシュのキャッシュが
    処理する．}
  \label{fig:l16-hybrid}
\end{marginfigure}

平均アクセス時間は，キャッシュの AMAT（\cref{eq:l12-amat}）と同じように
振る舞う．フラッシュから処理される要求はヒット時間を要し，ディスクが処理
する要求はミスの全体の時間を要する．コントローラは各ブロックがどこにあるか
を知っているので，ミスの際にまずフラッシュを読むことはなく，ミスペナルティ
は2つの時間の差である．
\caution{フラッシュはキャッシュであって，容量の追加ではない．8~GB の
フラッシュを持つ 2~TB のハイブリッドディスクは，2~TB のドライブとして
売られ，またそのように振る舞う．}

\begin{example}
  \label{ex:l16-hybrid}
  あるハイブリッドディスクが，アクセス時間 \SI{0.1}{\milli\second} の
  フラッシュと，回転を保ったままの \cref{ex:l16-access} のディスクとを
  組み合わせているとする．その例の要求について，ヒットでもコントローラと
  I/O バスの時間は支払うので $0.1 + 0.3 + 0.5 = \SI{0.9}{\milli\second}$
  であり，ミスは \SI{9}{\milli\second} かかるので，ミスペナルティは
  \SI{8.1}{\milli\second} である．要求の \SI{95}{\percent} がフラッシュで
  データを見つけるなら，平均アクセス時間は $0.9 + 0.05 \times 8.1 \approx
  \SI{1.3}{\milli\second}$ となり，フラッシュがデータのごく一部しか保持
  していないにもかかわらず，ディスク単体のおよそ7分の1で済む．いまやヒット
  の時間の大半は，コントローラと I/O バスが占める．\tip{ヒット時間は下限で
  ある．ヒット率をどれだけ高めても，このディスクはディスク単体の
  $9/0.9 = 10$ 倍を超えられ
  ない（Amdahl の法則，第2章）．}
\end{example}

\cref{tab:l16-compare} に，この節と前節の技術をまとめる．どれが望ましいか
は，何を最も重視するかによる．大量のデータを最も安く置くなら磁気ディスクが，
速度と低消費電力と頑健さならフラッシュが有利であり，アクセスがデータのごく
一部に集中するワークロードでは，ハイブリッドディスクが両方に近づく．

\begin{table}[htbp]
  \centering\small
  \caption{磁気ディスクと半導体ストレージの比較．}
  \label{tab:l16-compare}
  \begin{tabular}{@{}lllll@{}}
    \toprule
                      & 磁気ディスク & DRAM SSD     & フラッシュ SSD & ハイブリッド \\
    \midrule
    アクセス          & 遅い（ms）   & 最も速い     & 速い           & 大半は速い \\
    GB 当たりのコスト & 低い         & 非常に高い   & 中程度         & 低い \\
    消費電力          & 大きい       & 大きい       & 小さい         & より小さい \\
    可動部分          & あり         & なし         & なし           & あり \\
    無電源での保持    & する         & 電池が保つ間 & する           & する \\
    \bottomrule
  \end{tabular}
\end{table}

\begin{keypoint}
  ハイブリッドディスクは，キャッシュの原理をストレージに適用したもので
  ある．磁気ディスクの前にフラッシュを置くことで，大半の要求にフラッシュの
  速度と低い消費電力を与えつつ，ギガバイト当たりのコストはディスクのそれに
  ほぼ等しく保つ．
\end{keypoint}

\section{I/O デバイスの接続}
\label{sec:l16-io}

ストレージ装置は，他の入出力（I/O）デバイスと同様に，プロセッサに直接
取り付けられてはいない．\source{Lesson 16，動画 13--14．}業界標準で
定められたバスを通じて接続される．

\begin{definition}[I/O バス]
  \term[IO ばす]{I/O バス}[I/O bus]とは，I/O デバイスを計算機に接続する
  ための標準化されたバスである．標準はコネクタ，電気信号，プロトコルを
  定める．PCI Express，SATA，SCSI，USB がその例である．
\end{definition}

標準化によって，標準に従って作られたデバイスは，どちらを誰が製造したかに
よらず，そのバスを備えるどの計算機でも動作し，計算機の製造者は多様な
デバイスに対して1種類のインタフェースを用意すればよくなる．その代償は，
改善の速度が制限されることである．標準の新しい版は多くの企業の合意を必要と
し，古い版に合わせて作られたデバイスも動き続けるべきなので，バスはより速い
技術が利用可能になっても直ちにはそれを採用できない．ある標準に依存する
デバイスが多いほど，その標準の変化は遅くなる．\sidenote{PCI Express は
世代ごとに1レーン当たりの速度を倍にしてきた．2003 年の \num{2.5}~GT/s から
5，8，16（ここで7年待たされた），32，そして 2022 年の 64~GT/s である．
互換性は双方向に保たれ，
新しいカードは古いスロットでも古い速度で動く．}

I/O デバイスは，必要とするスループットも大きく異なる．グラフィックス
アダプタは膨大な量のデータを動かし，ディスクはそれよりはるかに少なく，
キーボードはほとんど動かさない．速いバス1本だけにすれば，すべてのデバイスが
高価で速いインタフェースを持たされることになり，遅いバス1本だけにすれば，
速いデバイスの性能が頭打ちになってしまう．そこでバスは階層をなすように配置
される（\cref{fig:l16-io-buses}）．

\begin{definition}[メザニンバス]
  \term[めざにんばす]{メザニンバス}[mezzanine bus]とは，プロセッサと，
  より遅い I/O バスとの間にある，速くて短い I/O バスである．今日ではそれは
  PCI Express である．グラフィックスアダプタのような高速なデバイスはこれに
  直接接続され，より遅いバスのコントローラもこれに取り付けられる．
\end{definition}

メザニンバス上の SATA コントローラや SCSI コントローラは，ディスクや光
ディスク装置のための SATA（あるいは SCSI）バスを提供し，USB ホスト
コントローラは，キーボードや USB メモリのような外部デバイスのための，ハブで
拡張される USB バスを提供する．\margin{SATA~3 は 6~Gb/s（600~MB/s），
USB 3.2 Gen~2 は 10~Gb/s，16 レーンの PCIe~5.0 スロットはおよそ 63~GB/s を
運ぶ．キーボードは毎秒数百バイトあればよい．}これらのバスは，本来的に遅く，
速度よりも
標準が安定していることの方が重要であるようなデバイスを扱う．メザニンバス
より遅く，変更される頻度も低い．PC のチップセットでは，SATA コントローラと
USB コントローラは1つのチップ，\term[さうすぶりっじ]{サウスブリッジ}%
[southbridge]に統合され，サウスブリッジはノースブリッジ（第15章）に接続
される．

\begin{figure}[htbp]
  \centering
  % Hierarchy of buses (logical view): processor and northbridge with main
% memory; the mezzanine bus (PCI Express) connects the graphics adapter
% directly and the SATA and USB controllers, which provide the slower buses
% (two point-to-point SATA links; a USB bus fanning out to its devices).
% In PC chipsets the two controllers form the southbridge (dashed box).
\begin{tikzpicture}[x=1cm,y=1cm, font=\footnotesize\sffamily, >={Stealth[length=1.6mm]},
  box/.style={draw, fill=ln-box, align=center, inner sep=3pt, minimum height=0.7cm},
  ctl/.style={draw, fill=ln-box, align=center, inner sep=2pt, minimum height=0.6cm, minimum width=2.1cm, font=\scriptsize\sffamily},
  dev/.style={draw, fill=white, align=center, inner sep=2pt, minimum height=0.6cm, font=\scriptsize\sffamily},
  lab/.style={font=\scriptsize\sffamily, text=ln-muted}]
  \node[box, minimum width=2.4cm] (cpu) at (5.0,3.3) {\iflangja{プロセッサ}{processor}};
  \node[box, minimum width=2.4cm, fill=ln-accent!20] (nb) at (5.0,2.05) {\iflangja{ノースブリッジ}{northbridge}};
  \node[dev, minimum width=1.9cm] (mem) at (0.95,2.05) {\iflangja{メインメモリ}{main memory}};
  \draw[<->, very thick, ln-accent] (cpu) -- node[lab, right] {FSB} (nb);
  \draw[<->, very thick, ln-accent] (nb) -- node[lab, above] {\iflangja{メモリバス}{memory bus}} (mem);
  % mezzanine bus
  \coordinate (bus) at (0,1.0);
  \draw[very thick, ln-accent] (nb.south) -- (nb.south |- bus);
  \draw[very thick, ln-accent] (1.25,1.0) -- (8.0,1.0);
  \node[lab, anchor=south west] at (5.12,1.02) {\iflangja{メザニンバス（PCI Express）}{mezzanine bus (PCI Express)}};
  \node[dev, minimum width=1.9cm] (gfx) at (1.25,0.0) {\iflangja{グラフィックス}{graphics}};
  \node[ctl] (sata) at (4.35,0.0) {\iflangja{SATA コントローラ}{SATA controller}};
  \node[ctl] (usb) at (8.0,0.0) {\iflangja{USB コントローラ}{USB controller}};
  \foreach \n in {gfx,sata,usb} {
    \draw[very thick, ln-accent] (\n.north) -- (\n.north |- bus);
  }
  % southbridge: the controllers of the slower buses
  \node[draw=ln-muted, dashed, rounded corners=2pt, fit=(sata)(usb), inner sep=4pt] (sb) {};
  \node[lab, anchor=west] at (sb.east) {\iflangja{サウスブリッジ}{southbridge}};
  % slower buses: two point-to-point SATA links, one USB bus
  \node[dev, minimum width=1.6cm] (d1) at (3.3,-1.55) {\iflangja{磁気ディスク}{magnetic disk}};
  \node[dev, minimum width=1.6cm] (d2) at (5.4,-1.55) {\iflangja{光ディスク装置}{optical drive}};
  \node[dev, minimum width=1.3cm] (d3) at (7.25,-1.55) {\iflangja{キーボード}{keyboard}};
  \node[dev, minimum width=1.3cm] (d4) at (8.75,-1.55) {\iflangja{USB メモリ}{flash drive}};
  \draw[thick] ([xshift=-3mm]sata.south) -- node[lab, left] {SATA} (d1.north);
  \draw[thick] ([xshift=3mm]sata.south) -- node[lab, right] {SATA} (d2.north);
  \coordinate (uj) at (8.0,-0.8);
  \draw[thick] (usb.south) -- (uj);
  \draw[thick] (uj) -- (d3.north);
  \draw[thick] (uj) -- (d4.north);
  \fill (uj) circle[radius=0.05];
  \node[lab, anchor=west] at (8.08,-0.6) {USB};
\end{tikzpicture}

  \caption{バスの階層．グラフィックスアダプタはメザニンバスに直接接続
    される．ディスクなどのより遅いデバイスには，メザニンバス上の SATA
    コントローラと USB コントローラを介して到達する．これらのコントローラ
    はより遅い標準 I/O バスを提供する．PC のチップセットでは，これらの
    コントローラがサウスブリッジ（破線）をなす．}
  \label{fig:l16-io-buses}
\end{figure}

各コントローラは，自身が提供するバスとメザニンバスとの間を橋渡しする．
ディスク読み出しのデータは，ディスクから SATA を通って SATA コントローラ
へ，メザニンバスを通り，ノースブリッジを抜けて主記憶へと進む．この経路が
\cref{eq:l16-access} の I/O バス時間である．速いデバイスはプロセッサの
近くに位置し，多数の遅いデバイスはコントローラの背後にある自分たちのバスに
閉じ込められて，速いデバイスの足を引っ張らない．\margin{新しいプロセッサ
はノースブリッジを統合し，サウスブリッジに直接接続する．}

技術が何であれ，どのように接続されていようと，ストレージはデータが生き残る
と期待される場所である．しかしあらゆるストレージ装置はいずれ故障する．
第17章では，そうした故障があってもシステムがデータとサービスを保ち続ける
方法を示す．

\begin{keypoint}[章のまとめ]
  ストレージはファイルと仮想記憶のページを永続的に保持し，そのレイテンシは
  DRAM のそれよりもさらにゆっくりとしか改善しない．磁気ディスクは，回転する
  プラッタのトラック上のセクタにデータを蓄える（\cref{eq:l16-capacity}）．
  そのアクセス（\cref{eq:l16-access}）はシーク時間と，平均で半回転である
  回転待ち時間とに支配され，容量が長年にわたり急速に伸びた一方で，この
  機械的なレイテンシはゆっくりとしか改善しない．光ディスクは可搬だが，
  標準化のために進歩が遅い．磁気テープは順次アクセスで安価な容量を提供し，
  バックアップに使われる．DRAM やフラッシュメモリで作られた SSD は，
  ギガバイト当たりの高いコストと引き換えに機械的な動作を避け，ハイブリッド
  ディスクはフラッシュを磁気ディスクのキャッシュとして使う．I/O デバイスは，
  階層をなすように配置された標準化された I/O バスを通じて接続される．速い
  デバイスは PCI Express のような速くて短いメザニンバスに直接つながり，
  より遅いデバイスは SATA や USB のような，より遅く安定したバスにつながって，
  それらのコントローラがメザニンバスに取り付けられる．
\end{keypoint}

\end{document}
